\section{Données du robot (baie CS9, serveur \texttt{uniVALplc})}
\label{secDonneesRobot}

La baie CS9 du robot SCARA est scrutée sur le réseau
\texttt{NOC\_ETHIP\_ROBOT}. Elle produit 148 octets et consomme 124 octets.
Contrairement au guichet, les échanges avec la baie CS9 ne se limitent pas à
des booléens élémentaires : ils sont structurés autour d'un objet unique,
de type \texttt{T\_StaeubliRobot}, utilisé comme paramètre d'entrée/sortie
par la quasi-totalité des blocs fonctionnels de la bibliothèque
\texttt{uniVALplc}.

\UPSTIremarque[Adresses mémoire non données]{
Les adresses \texttt{\%MW} de base des zones produites et consommées par
la baie CS9 ne sont pas données ici : leur détermination, à partir de la
cartographie mémoire du coupleur \texttt{NOC\_ETHIP\_ROBOT} (cf. document
\texttt{02c\_configuration\_noc}), fait partie du travail demandé aux
étudiants.}

\subsection{Principe des échanges}

À chaque cycle automate :
\begin{enumerate}
    \item les données produites par la baie CS9 (image mémoire de type
    \texttt{T\_FromRobot}) sont recopiées, via le bloc
    \texttt{VAL\_ReadAxesGroup}, dans les champs \texttt{Status} et
    \texttt{CommandInterface} de l'objet \texttt{T\_StaeubliRobot} ;
    \item le traitement utilisateur élabore, à partir de cet état, les
    ordres à transmettre au robot, en renseignant le champ
    \texttt{Command} de ce même objet ;
    \item ces ordres sont ensuite transcrits, via le bloc
    \texttt{VAL\_WriteAxesGroup}, dans l'image mémoire des données à
    consommer (type \texttt{T\_ToRobot}), transmise à la baie CS9.
\end{enumerate}

\subsection{Structure \texttt{T\_StaeubliRobot}}

\begin{center}
\begin{tabular}{|l|l|}
\hline
\textbf{Champ} & \textbf{Type} \\
\hline
\texttt{Status} & \texttt{T\_Status} (données produites -- lecture seule) \\
\hline
\texttt{Command} & \texttt{T\_Command} (données consommées -- écriture) \\
\hline
\texttt{CommandInterface} & \texttt{T\_CommandInterface} (interface de
communication) \\
\hline
\end{tabular}
\end{center}

\subsubsection{Champ \texttt{Status} (données produites -- extrait)}

\begin{center}
\begin{tabular}{|l|l|p{6.5cm}|}
\hline
\textbf{Champ} & \textbf{Type} & \textbf{Signification} \\
\hline
\texttt{Initialized} & BOOL & Vrai si la bibliothèque cliente est
initialisée \\
\hline
\texttt{Online} & BOOL & Vrai si le robot peut recevoir des commandes \\
\hline
\texttt{ErrorPending} & BOOL & Vrai si un défaut est en attente \\
\hline
\texttt{IsMoving} & BOOL & Vrai si le robot est en mouvement \\
\hline
\texttt{EStopActive} & BOOL & Vrai si un arrêt d'urgence est actif \\
\hline
\texttt{ActualSpeed} & REAL & Vitesse cartésienne courante du TCP \\
\hline
\texttt{ActualOperationMode} & UINT & Mode de fonctionnement courant
(0 = invalide, 1 = manuel, 3 = auto, 4 = distant) \\
\hline
\texttt{ActualErrorNumber} & UINT & Nombre de défauts en attente \\
\hline
\texttt{CartesianPos} & T\_CartesianPos & Position cartésienne courante \\
\hline
\texttt{JointPos} & T\_JointPos & Position articulaire courante \\
\hline
\texttt{RobotStateMachine} & UINT & État de la machine à états interne du
robot \\
\hline
\texttt{MovementID} & INT & Identifiant du mouvement en cours d'exécution \\
\hline
\texttt{MovementProgress} & INT & Pourcentage d'avancement du mouvement en
cours \\
\hline
\end{tabular}
\end{center}

\subsubsection{Champ \texttt{Command} (données consommées -- extrait)}

\begin{center}
\begin{tabular}{|l|l|p{6.5cm}|}
\hline
\textbf{Champ} & \textbf{Type} & \textbf{Signification} \\
\hline
\texttt{EnableVerbose} & BOOL & Active la trace d'activité du serveur \\
\hline
\texttt{OverrideCmd} & UINT & Consigne de vitesse (override), plage
$[1..100]$ \\
\hline
\texttt{OperationModeCmd} & UINT & Mode de fonctionnement demandé (0 =
invalide, 1 = manuel, 3 = auto, 4 = distant) \\
\hline
\texttt{ToolCmd} & UINT & Numéro de l'outil sélectionné pour le mouvement \\
\hline
\texttt{CoordSystemCmd} & UINT & Numéro du repère utilisateur sélectionné
pour le mouvement \\
\hline
\texttt{LifebitPeriod} & TIME & Période du bit de vie interne (entre
100~ms et 1000~ms) \\
\hline
\end{tabular}
\end{center}

\begin{UPSTIinfor}{Objet central des échanges}
L'objet \texttt{T\_StaeubliRobot} constitue le point de passage obligé de
tous les échanges avec la baie CS9 : c'est à la fois l'image de l'état du
robot (champ \texttt{Status}) et le support des ordres qui lui sont
adressés (champ \texttt{Command}).
\end{UPSTIinfor}
